\documentclass[10pt,a4paper]{amsart}
\usepackage{amsmath,amssymb,amsthm,graphicx}
\usepackage{verbatim}
\newenvironment{reference}{\comment}{\endcomment}
%\newenvironment{reference}{}{}
\newenvironment{slogan}{\comment}{\endcomment}
\newenvironment{history}{\comment}{\endcomment}

% For commutative diagrams we use Xy-pic
\usepackage[all]{xy}

% We use 2cell for 2-commutative diagrams.
\xyoption{2cell}
% Theorem environments.
%
\theoremstyle{plain}
\newtheorem{theorem}[subsection]{Theorem}
\newtheorem{proposition}[subsection]{Proposition}
\newtheorem{lemma}[subsection]{Lemma}

\theoremstyle{definition}
\newtheorem{definition}[subsection]{Definition}
\newtheorem{example}[subsection]{Example}
\newtheorem{exercise}[subsection]{Exercise}
\newtheorem{situation}[subsection]{Situation}

\theoremstyle{remark}
\newtheorem{remark}[subsection]{Remark}
\newtheorem{remarks}[subsection]{Remarks}

\numberwithin{equation}{subsection}

% Macros
%
\def\lim{\mathop{\mathrm{lim}}\nolimits}
\def\colim{\mathop{\mathrm{colim}}\nolimits}
\def\Spec{\mathop{\mathrm{Spec}}}
\def\Hom{\mathop{\mathrm{Hom}}\nolimits}
\def\Ext{\mathop{\mathrm{Ext}}\nolimits}
\def\SheafHom{\mathop{\mathcal{H}\!\mathit{om}}\nolimits}
\def\SheafExt{\mathop{\mathcal{E}\!\mathit{xt}}\nolimits}
\def\Sch{\mathit{Sch}}
\def\Mor{\mathop{\mathrm{Mor}}\nolimits}
\def\Ob{\mathop{\mathrm{Ob}}\nolimits}
\def\Sh{\mathop{\mathit{Sh}}\nolimits}
\def\NL{\mathop{N\!L}\nolimits}
\def\CH{\mathop{\mathrm{CH}}\nolimits}
\def\proetale{{pro\text{-}\acute{e}tale}}
\def\etale{{\acute{e}tale}}
\def\QCoh{\mathit{QCoh}}
\def\Ker{\mathop{\mathrm{Ker}}}
\def\Im{\mathop{\mathrm{Im}}}
\def\Coker{\mathop{\mathrm{Coker}}}
\def\Coim{\mathop{\mathrm{Coim}}}

% Boxtimes
%
\DeclareMathSymbol{\boxtimes}{\mathbin}{AMSa}{"02}

%
% Macros for moduli stacks/spaces
%
\def\QCohstack{\mathcal{QC}\!\mathit{oh}}
\def\Cohstack{\mathcal{C}\!\mathit{oh}}
\def\Spacesstack{\mathcal{S}\!\mathit{paces}}
\def\Quotfunctor{\mathrm{Quot}}
\def\Hilbfunctor{\mathrm{Hilb}}
\def\Curvesstack{\mathcal{C}\!\mathit{urves}}
\def\Polarizedstack{\mathcal{P}\!\mathit{olarized}}
\def\Complexesstack{\mathcal{C}\!\mathit{omplexes}}
% \Pic is the operator that assigns to X its picard group, usage \Pic(X)
% \Picardstack_{X/B} denotes the Picard stack of X over B
% \Picardfunctor_{X/B} denotes the Picard functor of X over B
\def\Pic{\mathop{\mathrm{Pic}}\nolimits}
\def\Picardstack{\mathcal{P}\!\mathit{ic}}
\def\Picardfunctor{\mathrm{Pic}}
\def\Deformationcategory{\mathcal{D}\!\mathit{ef}}
\usepackage{cleveref}
\crefname{originalreference}{original reference}{original references}
\Crefname{originalreference}{Original reference}{Original references}
\makeatletter
\newlabel{definition-weil-cohomology-theory-classical}{{0FGV}{source}{}{sourcepointer-1}{}}
\newlabel{definition-weil-cohomology-theory-classical@cref}{{[originalreference][0][]0FGV}{[1][0][]source}{}{}{}}
\makeatother
\crefname{originalreference}{original reference}{original references}
\Crefname{originalreference}{Original reference}{Original references}
\makeatletter
\newlabel{lemma-degrees-cycles-classical}{{0FGW}{source}{}{sourcepointer-2}{}}
\newlabel{lemma-degrees-cycles-classical@cref}{{[originalreference][0][]0FGW}{[1][0][]source}{}{}{}}
\makeatother
\crefname{originalreference}{original reference}{original references}
\Crefname{originalreference}{Original reference}{Original references}
\makeatletter
\newlabel{lemma-trace-product-classical}{{0FGX}{source}{}{sourcepointer-3}{}}
\newlabel{lemma-trace-product-classical@cref}{{[originalreference][0][]0FGX}{[1][0][]source}{}{}{}}
\makeatother
\crefname{originalreference}{original reference}{original references}
\Crefname{originalreference}{Original reference}{Original references}
\makeatletter
\newlabel{remark-replace-cup-product-classical}{{0FGU}{source}{}{sourcepointer-4}{}}
\newlabel{remark-replace-cup-product-classical@cref}{{[originalreference][0][]0FGU}{[1][0][]source}{}{}{}}
\makeatother
\title{Pairings of the Kunneth components of the class of the diagonal}
\author{The Stacks Project Authors}
\date{stacks-project:weil.tex:local-snapshot-20261012; GNU FDL 1.2 or later}
\begin{document}
\maketitle
\noindent\textit{Editorial note.} All delivered mathematical text is the author's own, copied from the pinned Stacks Project chapter file weil.tex; no statement or proof step was rewritten, reordered or removed.\par\medskip
\noindent\textit{Editorial note.} The counted proof is the author's complete argument: the identification of the class of the diagonal with the pushforward of the unit, the product formula for the degree map, and the choice of dual class producing the displayed pairings.\par\medskip
\noindent\textit{Editorial note.} The label lemma-class-diagonal-classical names the counted statement, and the proof environment immediately following it in the source is delivered verbatim as its proof.\par\medskip
\noindent\textit{Editorial note.} The counted statement is carried by the author's own lemma environment in the chapter source.\par\medskip
\noindent\textit{Editorial note.} Pre-open gate: python operations/agent-authoring/tools/precard\_gate.py --source weil.tex --label lemma-class-diagonal-classical --start 1752 --end 1799 --local-id 28142 returned RESULT=OK (L28142 free, occupancy FREE, public interval clear). Pre-assembly tag self-check: canonical tag resolved and not present in reference/public-items.jsonl.\par\medskip
\noindent\textit{Editorial note.} The statement and proof are copied from the pinned chapter file; the cited inputs are kept as external pointers with their canonical Stacks tags.\par\medskip
\noindent\textit{Allowed standard prerequisites (editorial).}\begin{itemize}
\item Pullback along the diagonal is the cup product. Same chapter, the classical version of the remark identifying the pullback along the diagonal with the cup product. Identifies the action of the diagonal on the cohomology of the product with the cup product.
\item The class of the diagonal as the pushforward of the unit. Same chapter, the classical version of the computation of the degree of a cycle, showing in particular that the class of the fundamental cycle is one and that classes of points do not vanish. Shows that the class of the diagonal is the pushforward of one, and that the class of the fundamental cycle is not zero.
\item Behaviour of the degree map under products. Same chapter, the classical version of the lemma computing the degree map of a product through the degree maps of the factors. Expands the pairing of the class of the diagonal with a product of classes into a sum of products of pairings on the factors.
\end{itemize}\medskip
\ifdefined\hypertarget\hypertarget{sourcepointer-1}{}\fi
\noindent\textit{Source reference.} Original reference 0FGV: sources/stacks-weil/pinned-original-source.tex, source lines 1682--1682. Prerequisite pointer invoked at proof line(s) 1755; it names a standard fact of the same chapter that the counted argument uses, and every step of the counted proof built on it is retained verbatim in the delivered body.\par\medskip
\ifdefined\hypertarget\hypertarget{sourcepointer-2}{}\fi
\noindent\textit{Source reference.} Original reference 0FGW: sources/stacks-weil/pinned-original-source.tex, source lines 1695--1695. Prerequisite pointer invoked at proof line(s) 1781; it names a standard fact of the same chapter that the counted argument uses, and every step of the counted proof built on it is retained verbatim in the delivered body.\par\medskip
\ifdefined\hypertarget\hypertarget{sourcepointer-3}{}\fi
\noindent\textit{Source reference.} Original reference 0FGX: sources/stacks-weil/pinned-original-source.tex, source lines 1719--1719. Prerequisite pointer invoked at proof line(s) 1793; it names a standard fact of the same chapter that the counted argument uses, and every step of the counted proof built on it is retained verbatim in the delivered body.\par\medskip
\ifdefined\hypertarget\hypertarget{sourcepointer-4}{}\fi
\noindent\textit{Source reference.} Original reference 0FGU: sources/stacks-weil/pinned-original-source.tex, source lines 1653--1653. Prerequisite pointer invoked at proof line(s) 1773; it names a standard fact of the same chapter that the counted argument uses, and every step of the counted proof built on it is retained verbatim in the delivered body.\par\medskip
\begin{lemma}
\label{lemma-class-diagonal-classical}
Let $H^*$ be a classical Weil cohomology theory
(Definition \ref{definition-weil-cohomology-theory-classical}).
Let $X$ be a smooth projective variety of dimension $d$.
Choose a basis $e_{i, j}, j = 1, \ldots, \beta_i$ of $H^i(X)$ over $F$.
Using K\"unneth write
$$
\gamma([\Delta]) =
\sum\nolimits_{i = 0, \ldots, 2d}
\sum\nolimits_j e_{i, j} \otimes e'_{2d - i , j}
\quad\text{in}\quad
\bigoplus\nolimits_i H^i(X) \otimes_F H^{2d - i}(X)
$$
with $e'_{2d - i, j} \in H^{2d - i}(X)$.
Then $\int_X e_{i, j} \cup e'_{2d - i, j'} = (-1)^i\delta_{jj'}$.
\end{lemma}

\begin{proof}
Recall that $\Delta^* : H^*(X \times X) \to H^*(X)$ is equal to the
cup product map $H^*(X) \otimes_F H^*(X) \to H^*(X)$, see
Remark \ref{remark-replace-cup-product-classical}. On the other hand we have
$\gamma([\Delta]) = \Delta_*\gamma([X]) = \Delta_*1$ by
axiom (C)(b) and the fact that $\gamma([X]) = 1$. Namely,
$[X] \cdot [X] = [X]$ hence by axiom (C)(c) the cohomology class
$\gamma([X])$ is $0$ or $1$ in the $1$-dimensional $F$-algebra $H^0(X)$;
here we have also used axioms (A)(d) and (A)(b).
But $\gamma([X])$ cannot be zero as $[X] \cdot [x] = [x]$
for a closed point $x$ of $X$ and we have the nonvanishing
of $\gamma([x])$ by Lemma \ref{lemma-degrees-cycles-classical}.
Hence
$$
\int_{X \times X} \gamma([\Delta]) \cup a \otimes b =
\int_{X \times X} \Delta_*1 \cup a \otimes b =
\int_X a \cup b
$$
by the definition of $\Delta_*$. On the other hand, we have
$$
\int_{X \times X} (\sum e_{i, j} \otimes e'_{2d -i , j}) \cup a \otimes b =
\sum (\int_X a \cup e_{i, j})(\int_X e'_{2d - i, j} \cup b)
$$
by Lemma \ref{lemma-trace-product-classical}; note that we made
two switches of order so that the sign is $1$.
Thus if we choose $a$ such that $\int_X a \cup e_{i, j} = 1$
and all other pairings equal to zero, then we conclude that
$\int_X e'_{2d - i, j} \cup b = \int_X a \cup b$ for all $b$, i.e.,
$e'_{2d - i, j} = a$. This proves the lemma.
\end{proof}

\end{document}
